\subsection{محدودیت‌ها و تحلیل شکست}
تحلیل رخدادهای بنچ‌مارک و بررسی کد منبع موارد زیر را نمایان می‌سازد:
\begin{itemize}
    \item \textbf{رخداد خطای زمان اجرا در بنچ‌مارک دیکودری \lr{GPT-2}:} در ارزیابی مدل \lr{GPT-2} بر روی مجموعه دادگان WikiText-2، بنچ‌مارک دچار وضعیت کرش گردید (\lr{CRASH}). بررسی ردپای خطا حاکی از بروز استثنای متغیر محلی در مرحله ارزیابی دقت است (\lr{UnboundLocalError: cannot access local variable 'majority\_count'}). این باگ نرم‌افزاری ناشی از شرط‌بندی نادرست بر روی ابعاد لاجیت‌های توکن‌های دنباله دیکودری در اسکریپت درایور بنچ‌مارک بوده و ناشی از ماهیت ریاضی خود کوانتیزه‌ساز نیست.
    \item \textbf{پیچیدگی شتاب‌بخشی سخت‌افزاری در پردازش بیت‌های ناهمگن:} اگرچه کوانتیزه‌سازی شبیه‌سازی‌شده سربار حافظه تئوریک را کاهش می‌دهد، اما اجرای عملیاتی کانال‌های با طول بیت‌های متفاوت (مثلاً ۲ بیت در کنار ۸ بیت) روی کرنل‌های \lr{GEMM} یکپارچه پردازنده‌های گرافیکی نیازمند کتابخانه‌های تخصصی چیدمان بیت (\lr{Bit-Packing}) است.
\end{itemize}
